\chapter{QM/MM input files (advanced)}
\label{ch:qmmm}

Some questions about a redox site, such as where an electron sits or how the
site's energy changes, need quantum chemistry. A whole protein is far too big
for that. In a QM/MM calculation only the site is treated with quantum
chemistry, and the rest of the protein with the ordinary force field. The
\menupath{5c}{QM/MM Preparator} writes the input files for such a
calculation. It does not run the calculation; that needs Gaussian or ORCA,
which you install and license yourself.

\begin{glance}
\glanceitem{What it does}{Writes a Gaussian ONIOM input file (\file{.com})
  for a redox site inside its protein, with layers, link atoms, frozen atoms
  and all Amber parameters. It can also prepare ORCA inputs if ORCA is
  installed.}
\glanceitem{You need it when}{You want to run a QM/MM calculation on a heme,
  flavin, iron--sulfur cluster or other redox site.}
\glanceitem{Before you start}{Redox sites from 2a (Chapter~\ref{ch:redox}),
  usually transformed with 4d, and a topology and coordinates
  (\file{.prmtop} and \file{.rst7}) from 5a (Chapter~\ref{ch:topology}).
  Gaussian~16 to run the file. ORCA, including \texttt{orca\_mm}, even to
  prepare ORCA inputs.}
\glanceitem{Where to find it}{\key{5c} QM/MM Preparator in the main menu.}
\glanceitem{You get}{\file{oniom_inputs/oniom_1.com} plus a plain-text
  summary (\file{oniom_summary_1.txt}).}
\end{glance}

\section{QM/MM in plain words}
\label{sec:qmmm-words}

\term{QM} (quantum mechanics) calculates the electrons explicitly. It can
describe bonds breaking, charges moving and different spin states, but it is
slow: a few hundred atoms at most. \term{MM} (molecular mechanics) is the
force field used for MD. It is fast but cannot describe electrons. A
\term{QM/MM} calculation treats a small region with QM and the rest with MM.

Gaussian's way of doing this is called \term{ONIOM}. It splits the system into
\term{layers}:

\begin{itemize}
  \item the \textbf{HIGH} layer: the QM region, for example the heme and the
        side chains bound to the iron;
  \item an optional \textbf{MEDIUM} layer, treated with a cheaper QM method
        (three-layer ONIOM only);
  \item the \textbf{LOW} layer: everything else, treated with the Amber force
        field.
\end{itemize}

Where a bond crosses from the QM region to the MM region, the QM side would
have a dangling bond. ONIOM caps it with a hydrogen called a
\term{link atom}. In \term{electronic embedding} the MM partial charges act on
the QM electrons; in \term{mechanical embedding} they do not. MM charges very
close to the QM region can pull the electrons too hard, so ProPrep offers ways
to reduce them (Section~\ref{sec:qmmm-boundary}).

Two numbers must be right for the QM region: its \term{charge} and its
\term{multiplicity}. The multiplicity is the number of unpaired electrons plus
one: 1 (singlet) if all electrons are paired, 2 (doublet) for one unpaired
electron, and so on. Low-spin Fe(III), as in ferric cytochrome~c, has one
unpaired electron, so its multiplicity is 2.

\section{Step by step: Gaussian ONIOM for cytochrome c}
\label{sec:qmmm-oniom}

This continues the \pdb{1HRC} example of Chapter~\ref{ch:redox}. The heme site
was detected in 2a, transformed in 4d (ferric, fixed pH, both propionates
deprotonated), and built in 5a into a topology and coordinate file
(\file{.prmtop} and \file{.rst7}; ff14SB, implicit solvent, total charge
$+7$).

\subsection*{Choose the structure and the program}

\begin{steps}
  \item Type \key{5c}, then \key{1} (\prompt{Run QM/MM preparation
        workflow}).
  \item ProPrep finds the topology and coordinates from 5a:
\begin{screen}
  Found topology in workspace: final_transformed_structure.prmtop
Use this topology? (Enter to accept, or 'b' to browse) ():
  Prmtop: final_transformed_structure.prmtop
Coordinate source: [r]st7/inpcrd or [t]rajectory (r): §ans«r»
  Found coordinates in workspace: final_transformed_structure.rst7
Use these coordinates? (Enter to accept, or 'b' to browse) ():
...
Select QM/MM interface:
  1. Gaussian (ONIOM)
  2. ORCA
Select interface (1): §ans«1»
\end{screen}
        Press \key{Enter} twice to accept the files, and choose \key{1}.
\end{steps}

The ONIOM checklist appears. Its state is saved in
\file{oniom_inputs/workflow_state.json} after every step.

\begin{screen}
  SETUP
  [ ]  1.  Select RedoxSite                     Choose which redox site(s) ...
  [ ]  2.  Configure ONIOM Layers               Assign residues to HIGH/MED...
  CONFIGURATION
  [ ]  3.  Configure MM Settings & Atom Typing  Select force field, run ato...
  [ ]  4.  Configure QM Settings                Set QM functional, basis se...
  [ ]  5.  Configure Job Settings (optional)    Set job type, processors, m...
  GENERATION
  [ ]  6.  Prepare & Validate ONIOM Setup       Build ONIOM configuration, ...
  [ ]  7.  Write Gaussian Input File            Write production or diagnos...
...
  [n]ext (1. Select RedoxSite)  |  [v]iew structure  |  [s]ave & quit
  [q]uit  |  [#] to run/re-run step  |  [#s] skip optional step
\end{screen}

Type \key{n} to go to the next step. A number re-runs a step; re-running a
step resets the later ones to ``pending''.

\subsection*{Step 1: select the site}

At \prompt{Select RedoxSite(s) for ONIOM calculation (e.g. 1, 1-3, 1,3,5, or
all) (1)} type \key{1}. Several sites would be merged into one QM region.
ProPrep confirms \prompt{Selected RedoxSite 1 with 73 atoms}.

\subsection*{Step 2: the layers}

\begin{steps}
  \item \prompt{Number of ONIOM layers (2 or 3) (2)}: \key{2}.
  \item ProPrep puts the site's residues into the HIGH layer.
        \textbf{Read this list carefully}:
\begin{screen}
✓ Automatically added 7 residues from RedoxSite to HIGH layer:
    CYO14 (6 atoms)
    CYO17 (6 atoms)
    HIO18 (6 atoms)
    MEO80 (6 atoms)
    HMO105 (87 atoms)
    PRD107 (9 atoms)
    WAT109 (3 atoms)
\end{screen}
        For this site the list should hold the heme (HMO), \emph{both}
        propionates (two PRD) and the four backbone stubs (CYO, CYO, HIO,
        MEO), and no water. Here a water (WAT109) was added and one
        propionate (PRD106) was missing.
  \item Fix it. Answer \key{y} to \prompt{Add or remove whole residues?},
        type \typed{-109} to remove the water and \typed{106} to add the
        propionate, then press \key{Enter} on an empty line:
\begin{screen}
  HIGH residue number (or Enter to finish) (): §ans«-109»
✓ Removed WAT109
  HIGH residue number (or Enter to finish) (): §ans«106»
✓ Added PRD106 (9 atoms)
  HIGH residue number (or Enter to finish) ():

HIGH layer: 7 residues
\end{screen}
  \item \prompt{Trim any residues to sidechain only (CA-CB cut)? (n)}: press
        \key{Enter} (or answer \key{y} and then press \key{Enter} to skip).
  \item \prompt{Freeze atoms beyond distance from QM region (Å, 0 for none)
        (0)}: we typed \typed{12}. Atoms more than 12~\AA{} from the QM
        region will not move in an optimization.
\end{steps}

\begin{watchout}
\textbf{Always check the HIGH-layer list by residue name.} In version 1.23.0,
5c matches site residues to topology residues by number. In 1HRC, 4d had
written the propionates between crystal waters, so the numbers were shifted:
ONIOM silently put a water into the QM region and left a propionate out. The
calculation would run, but on the wrong molecule. The numbers to type are the
ones shown in this list. In our run the heme was HMO105, so its propionates
were PRD106 and PRD107.
\end{watchout}

\begin{watchout}
The trim question in the same step numbers residues differently: it listed
the propionates as \texttt{[1] PRD107} and \texttt{[2] PRD108}, one higher than
the add/remove prompt. It also offers the propionates for trimming because
their atoms are called CA and CB; trimming them would be chemically wrong.
Do not trim propionates.
\end{watchout}

\subsection*{Step 3: MM settings and atom typing}

ProPrep reads atom types and charges from the topology (\prompt{Extracted 2121
atoms, 39 unique atom types}). It finds the bonds that cross the QM/MM
boundary (here 6, at the C$\alpha$ atoms next to the four stub residues) and
adds small caps from the neighboring residues to the HIGH layer
(\prompt{Expanded HIGH by 12 atom(s)}). It then prints the charge of every
HIGH residue (total $-1.1126$), a summary of the LOW layer ($+8.1126$), and
writes \file{debug_atom_charges.txt}.

\label{sec:qmmm-boundary}
Next it asks how to treat the MM charges right at the boundary:

\begin{screen}
Boundary charge treatment options:
  1. No adjustment
  2. Gaussian ScaleCharge keyword (recommended for electronic embedding)
  3. Charge-shift redistribution
  4. Manual per-atom override
Select boundary charge treatment (1): §ans«2»
...
Enter ScaleCharge value (digits only, e.g. 500) (500):
\end{screen}

For electronic embedding choose \key{2} and press \key{Enter} to accept
\texttt{500}.

\begin{itemize}
  \item \textbf{1 No adjustment}: keep all charges. For mechanical embedding.
  \item \textbf{2 ScaleCharge}: Gaussian scales charges near the QM region
        while it runs. Each digit, read from the right, is the scale for
        atoms 1, 2, 3 \ldots{} bonds away, in steps of 0.2 (0 = charge off,
        5 = full charge). \typed{500} (Gaussian's default) switches off
        charges within two bonds. The digits must not increase from left to
        right.
  \item \textbf{3 Charge shift}: sets the C$\alpha$ charges at the boundary to
        zero and moves their charge to the neighboring atoms, keeping the
        total.
  \item \textbf{4 Manual}: type a new charge for each boundary atom.
\end{itemize}

\begin{watchout}
ProPrep does not check the ScaleCharge digits. In our test \typed{135}
(digits increasing, against the rule printed on screen) was accepted. Also,
if you re-run step 3 and choose another treatment, the earlier
\texttt{ScaleCharge=} stays in the route line, so both are applied. Check the
route line in step 5's preview, and if needed delete the old keyword from the
first lines of the \file{.com} file with a text editor.
\end{watchout}

Step 3 ends with a summary that includes
\prompt{Suggested QM region charge: -1} and \prompt{Total system (real)
charge: 7}. Answer \key{n} to \prompt{View ONIOM layers in 3D structure
viewer?} unless you want the browser viewer (not tried for this manual).

\subsection*{Step 4: QM settings}

The first time, we typed \typed{B3LYP} and \typed{def2-SVP} at the first two
questions (the basis-set default is \texttt{6-31G*}). This screen is from our
second pass through step 4, where those answers had become the defaults:

\begin{screen}
QM functional (e.g., B3LYP, M06-2X, wB97X-D) (B3LYP):
QM basis set (e.g., 6-31G*, 6-311G(d,p), def2-SVP) (def2-SVP):
QM region charge (HIGH layer total: -1.11) (-1):
Multiplicity Guide:
  1 = singlet (paired electrons)
  2 = doublet (1 unpaired electron)
  3 = triplet (2 unpaired electrons)
QM region multiplicity (1): §ans«2»
\end{screen}

The functional and basis set are passed straight to Gaussian. The QM charge
of $-1$ comes from Fe(III) heme ($+1$) and two charged propionates ($-2$).
The multiplicity is \textbf{2} for low-spin Fe(III).

\begin{watchout}
The multiplicity default is 1, and ProPrep does not check it against the
number of electrons. In our first run we accepted 1 for the ferric heme.
ProPrep then wrote the charge/multiplicity line \texttt{7 2 -1 1}: doublet for
the whole system, singlet for the QM region. Gaussian rejects such a pair,
because both regions have an odd number of electrons. To fix it, re-run step
4 (type \key{4}, then \key{y}) with multiplicity 2, then run steps 5--7 again.
\end{watchout}

A yellow \prompt{Open-shell / BS-DFT advisory} always follows. It reminds you
that metal clusters with coupled spins (such as [4Fe-4S]) need a special
multi-step procedure in Gaussian (\url{https://gaussian.com/afc/}), and that
ProPrep does not set this up.

\subsection*{Step 5: job settings (optional)}

\begin{screen}
Job type [Opt/SP/Freq/Opt Freq] (Opt): §ans«SP»
Number of processors (4): §ans«8»
Memory (GB) (8): §ans«16»
Route line preview:
  #P ONIOM(B3LYP/def2-SVP:AMBER=SoftOnly)=(EmbedCharge,ScaleCharge=135) SP ...
Enter any extra Gaussian keywords to append, or press Enter to skip.
Extra keywords (): §ans«SCF=XQC»
\end{screen}

In our run the route shows \texttt{ScaleCharge=135}, the unchecked value
described in the watch-out box above; with the default you would see
\texttt{ScaleCharge=500}.

\texttt{SP} is a single-point energy; \texttt{Opt} a geometry optimization;
\texttt{Freq} vibrational frequencies. Extra keywords are added to the route
line as typed. Type \typed{5s} instead of \key{n} to skip this step and keep
the defaults.

\subsection*{Step 6: prepare and validate}

Nothing to answer. ProPrep freezes distant atoms (\prompt{Froze 359 atoms
beyond 12.0Å from any HIGH atom}), places link atoms (\prompt{Placed 6 link
atoms}), builds the bond list, copies all force-field parameters from the
topology into the file, and checks the set-up:

\begin{screen}
  Total atoms: 2121
  HIGH: 141, MEDIUM: 0, LOW: 1980
  Link atoms: 6
  Frozen: 359, Active: 1762
  Total system charge: 7.000003 (nearest integer: 7)
✓ Validation passed
\end{screen}

\subsection*{Step 7: write the file}

\begin{screen}
Write Gaussian Input
1. Write production input file
2. Write diagnostic file (ONIOM=OnlyInputFiles)
   Gaussian will write separate .com files for each sub-calculation
   so you can verify layer assignments. No calculation is run.
Select option (1): §ans«1»
Output filename (oniom_1.com):
...
  ✓ Successfully wrote ONIOM input file (199.0 KB)
...
Summary report filename (oniom_summary_1.txt):
\end{screen}

ProPrep writes the input file \file{oniom_1.com} and a summary into the
folder \file{oniom_inputs}. It also writes \file{oniom_1.txt}, which has the same
content as the summary.

\section{Checking the files before you run Gaussian}
\label{sec:qmmm-check}

Open both files in a text editor. Start with the summary
(\file{oniom_summary_1.txt}):

\begin{screen}
Layer distribution:
  HIGH:    141 atoms in  13 residues (charge: -1.11)
  LOW:    1980 atoms in 225 residues
  LINK:      6 hydrogen link atoms
...
QM Method:
  Functional: B3LYP
  Basis set:  def2-SVP
  Charge:     -1
  Multiplicity: 2
...
HIGH layer residues:
  LYS14
  CYO15
...
  HMO106
  PRD107
  PRD108
\end{screen}

The HIGH residues include the small caps from the neighboring residues (LYS,
ALA, GLN, THR, ILE). Check by \emph{name} that HMO, both PRD and the four
stubs are there and no WAT. The numbers here are one higher than at the step-2
prompt.

Then look at the first lines of the \file{.com} file. This is the file from
our first pass, which still had both mistakes described above (ScaleCharge
135 and QM multiplicity 1):

\begin{screen}
%mem=16GB
%nprocshared=8
#P ONIOM(B3LYP/def2-SVP:AMBER=SoftOnly)=(EmbedCharge,ScaleCharge=135) SP ...

ONIOM 2-layer calculation: HIGH=141 atoms, MEDIUM=0 atoms, LOW=1980 atoms, ...

7 2 -1 1
H-HC-0.112300        -1    49.817000    10.944000    -5.932000 L
C-CT--0.366200       -1    50.594000    11.606000    -6.314000 L
...
C-CX-0.000000         0    47.451000    29.679000    -2.940000 L   H-HX-0.0  222
\end{screen}

\begin{itemize}
  \item \texttt{\%mem} and \texttt{\%nprocshared} are the memory and
        processors from step 5.
  \item The \texttt{\#P} line is the \term{route line}: the method for each
        layer and the options. Check that only one boundary treatment is in
        it.
  \item The \textbf{charge/multiplicity line}: the first pair is the whole
        system (charge $+7$, doublet), the second pair is the QM region
        (charge $-1$). Both multiplicities must fit the electron count; for
        this ferric heme both must be 2. After we re-ran step 4 with
        multiplicity 2, the line read \texttt{7 2 -1 2}.
  \item Each atom line has element--type--charge, a freeze flag (0 = moves,
        $-1$ = frozen), the coordinates and the layer (H, M or L). A line
        ending in \texttt{H-HX-0.0 222} is a boundary atom with its link
        hydrogen.
  \item After the atoms come the bond list and the Amber parameters
        (\texttt{AmbTrs}, \texttt{HrmBnd1}, \texttt{VDW}, \texttt{NonBon}).
\end{itemize}

\begin{tip}
Write a diagnostic file too: step 7, option \key{2}. It writes
\file{oniom_1_diagnostic.com} with \texttt{OnlyInputFiles} in the route line.
Gaussian then only writes out the separate sub-calculations without running
them, so you can check the layers cheaply. ProPrep's message mentions a
``Test keyword'', but the file contains \texttt{OnlyInputFiles}, not
\texttt{Test}.
\end{tip}

\section{What you do next with Gaussian}
\label{sec:qmmm-gaussian}

\begin{needsoftware}
Running the \file{.com} file needs \textbf{Gaussian~16}, which is licensed
separately. ProPrep does not submit or run it. Gaussian was not available
while this manual was written, so the files above were not run.
\end{needsoftware}

On a computer with Gaussian, go to the \file{oniom_inputs} folder and run:

\begin{terminal}
g16 < oniom_1.com > oniom_1.log
\end{terminal}

Make sure the computer has at least the memory and processors written at the
top of the file. On a computing cluster, use the cluster's job system
instead; ask your system administrator. When the job finishes, the end of
\file{oniom_1.log} should say ``Normal termination''. You can open the log
in a viewer such as GaussView. Interpreting QM/MM results is outside the scope
of this manual.

\section{Other choices}
\label{sec:qmmm-other}

\subsection{Three-layer ONIOM}

At \prompt{Number of ONIOM layers} type \key{3}. After the HIGH layer
ProPrep asks for MEDIUM residues (we typed \typed{28} and \typed{58}, two
threonines). There is no question for the MEDIUM method: it is fixed at
HF/3-21G, giving the route
\texttt{ONIOM(B3LYP/def2-SVP:HF/3-21G:AMBER=SoftOnly)}. The charge line has
four pairs; in our test it read \texttt{7 2 0 1 0 1 0 2}. Have someone
experienced in ONIOM check these values.

\begin{watchout}
In our three-layer test we kept the faulty automatic HIGH list (water in,
PRD106 out). The heme--PRD106 bond then crossed the QM/MM boundary, but
ProPrep did not detect it: no link atom, no warning. Fix the HIGH list first.
\end{watchout}

\subsection{Manual boundary charges}

Option \key{4} asks for a new charge for each boundary atom. Typing \typed{x}
keeps the original. If the total no longer adds up to a whole number, ProPrep
warns \prompt{Charge may be unbalanced} and later
\prompt{Total charge (7.2400) deviates from nearest integer (7)}, but still
passes validation. Keep the total charge whole.

\subsection{Coordinates from an MD trajectory}

At \prompt{Coordinate source} type \key{t} to take frames from a trajectory
(Chapter~\ref{ch:runningmd}). ProPrep asks for the trajectory file(s),
\prompt{Center solute in box (autoimage)?}, and a frame range
\texttt{start:end:increment} (shown on screen with an emoji in place of
``:end:''). Extracting one frame, e.g. \typed{4:4:1}, worked. Extracting
several frames failed in version 1.23.0
(\prompt{cpptraj produced no output frames}), and the single-frame run that
followed it picked up a left-over PDB from the failed run. Extract one frame
at a time, in a fresh session.

\subsection{Save, resume and reset}

\key{s} saves and leaves (\prompt{State saved. You can resume later.}). The
next time you choose Gaussian, a \prompt{Resume Available} panel offers to
continue (\prompt{Progress: 6/7 steps completed}).

\begin{watchout}
Answering \key{n} to \prompt{Resume from saved state?} does not fully start
fresh: in our test the earlier functional, basis set, job type, processors,
memory, ScaleCharge and extra keywords all reappeared in the next file. Check
every setting again, or use 5c option \key{3} (\prompt{Reset configuration to
defaults}) first.
\end{watchout}

Menu option \key{2} shows the current configuration. Before any run it says
\prompt{No QM/MM workflow has been run yet.}

\subsection{ORCA}
\label{sec:qmmm-orca}

\begin{needsoftware}
The ORCA route needs ORCA's \texttt{orca\_mm} program \emph{while preparing}
the files, not only to run them: step 2 converts the Amber topology into
ORCA's force-field format with \texttt{orca\_mm -convff -AMBER}. Without it,
ProPrep shows \prompt{orca\_mm not found on PATH}, step 2 fails, and the last
step fails too (\prompt{Step 8 failed: [Errno 2] No such file or directory:
''}). \textbf{No ORCA files are written.}
\end{needsoftware}

Without ORCA we could still walk through the questions. The ORCA checklist
has eight steps: verify topology, convert force field, select QM region,
active region (atoms within a distance, default 6~\AA, may move), QM settings
(method, basis, dispersion correction D3BJ, charge, multiplicity), QM/MM
coupling (electrostatic or mechanical embedding, charge scheme CS, RCD or
Z0--Z3, link-atom distances), job settings, and preview \& write.

Two things are better than in the ONIOM route. In step 3 ProPrep warns
\prompt{5 atom(s) could not be mapped} for the misplaced propionate, and for
the bond between the heme and PRD106 it offers
\prompt{[1] Include entire PRD106 in QM (9 atoms)}, which fixes the problem.
On the other hand, typing a non-number at a number question (for example
\typed{abc} for the active-region distance) makes the step fail with a Python
error message; run the step again.

According to ProPrep's source, with \texttt{orca\_mm} present it would write
an ORCA \file{.inp} file, a matching PDB, the converted \file{.prms}
force-field file and a run script in \file{orca_inputs/}. We could not check
these files.

\section{If something goes wrong}
\label{sec:qmmm-problems}

\begin{itemize}\sloppy
  \item \textbf{\prompt{Frame extraction failed: No prmtop file selected}.}
    There is no topology yet. Build one with 5a first
    (Chapter~\ref{ch:topology}). The main menu only says that 5c needs redox
    sites.
  \item \textbf{A water is in the HIGH layer, or a propionate is missing.}
    Remove and add residues in step 2 (Section~\ref{sec:qmmm-oniom}).
  \item \textbf{The charge/multiplicity line looks inconsistent.} Re-run step
    4 with the right multiplicity, then steps 5--7.
  \item \textbf{Two boundary treatments in the route line.} Delete the old
    \texttt{ScaleCharge=} from the \file{.com} file.
  \item \textbf{Old settings come back after declining Resume.} Use 5c option
    3 (reset) before starting again.
  \item \textbf{ORCA step 2 fails.} Install ORCA and give the folder that
    contains \texttt{orca\_mm}; there is no way round it.
  \item \textbf{Several trajectory frames fail.} Extract one frame at a time.
\end{itemize}

See also Chapter~\ref{ch:troubleshooting} and Chapter~\ref{ch:knownissues}.
